\documentclass[12pt]{article}
\usepackage[utf8]{inputenc}
\usepackage[T1]{fontenc}
\usepackage{hyperref}
\title{Soyaltepec Mazatec Language Resources for K-12 Curriculum Development}
\author{Compiled by Claude Sonnet 5 (claude-sonnet-5) for Andr\'as Kornai}
\date{}
\begin{document}
\maketitle
\section{Introduction}
Soyaltepec Mazatec (ISO 639-3 vmp, Glottocode soya1237) is an Otomanguean language of Oaxaca, Mexico, around San Felipe Jalapa de D\'iaz. The search results disagree on its vitality. Wikitongues gives EGIDS 6a (vigorous), while Wikipedia gives 28,000 speakers (2000) but also says perhaps only about 900 people, mostly monolingual, speak the original variety. The search found no school materials, no dictionary and no modern curriculum resources specific to this variety.
\section{Primary Educational Resources}
The Colegio de M\'exico repository holds 1954 primers (cartillas) in Mazateco de Soyaltepec, published with the Instituto Ling\"u\'istico de Verano. The repository notes that the language of this edition is no longer current but is useful as a precursor to later editions.
\begin{itemize}
\item Colegio de M\'exico repository, primer record 1: \url{https://repositorio.colmex.mx/concern/books/d791sj007?locale=en}
\item Colegio de M\'exico repository, primer record 2: \url{https://repositorio.colmex.mx/concern/books/dz010r92r?locale=en}
\end{itemize}
\section{Linguistic and Reference Materials}
\begin{itemize}
\item Wikipedia article, which cites Heather Beal's 2011 thesis on the segments and tones of Soyaltepec Mazatec: \url{https://en.wikipedia.org/wiki/Soyaltepec_Mazatec}
\item Wikitongues language page: \url{https://wikitongues.org/languages/vmp/}
\item ASJP wordlist (San Miguel Soyoltepec): \url{https://asjp.clld.org/languages/SAN_MIGUEL_SOYOLTEPEC}
\end{itemize}
No dictionary or full grammar turned up in the search results.
\section{Religious and Community Resources}
\begin{itemize}
\item Scripture Earth listing, which points to portions of scripture on YouVersion: \url{https://scriptureearth.org/00eng.php?iso=vmp}
\item Joshua Project language page: \url{https://joshuaproject.net/index.php/languages/vmp}
\end{itemize}
\section{Digital and Audio Resources}
\begin{itemize}
\item Global Recordings Network, audio recordings in Mazateco, Soyaltepec: \url{https://www.globalrecordings.net/en/language/vmp}
\end{itemize}
\section{Curriculum Development Context}
Little usable curriculum material was found. A curriculum would have to start from the 1954 primers, the audio recordings and the scripture portions, together with new documentation work with the community. The conflicting vitality estimates should be settled by local fieldwork before planning.
\end{document}